% Figure: Nusselt number for natural convection in a vertical enclosed air
% gap versus the gap Rayleigh number. Below a critical Rayleigh number the
% gap conducts (Nu = 1, still layer); above it convection cells form and
% Nu climbs as a power law. Log-log axes, schematic of the two regimes.
\begin{figure}[htbp]
\centering
\begin{tikzpicture}
\begin{axis}[
  width=12cm, height=7cm,
  xmode=log, ymode=log,
  xlabel={gap Rayleigh number $\text{Ra}_L$},
  ylabel={gap Nusselt number $\text{Nu}$},
  xmin=10, xmax=1e6,
  ymin=0.8, ymax=20,
  axis lines=left,
  tick label style={font=\scriptsize},
  label style={font=\small},
  legend style={font=\scriptsize, at={(0.02,0.98)}, anchor=north west},
  legend cell align=left,
  legend entries={conduction floor $\text{Nu}=1$,convection branch $0.197\,\text{Ra}^{1/4}$},
  clip=false,
]
% conduction plateau Nu = 1 up to Ra ~ 2000
\addplot[engblue, very thick, domain=10:2000, samples=20] {1};
% convection branch Nu = 0.197 Ra^{1/4} (a representative vertical-cavity fit)
% at Ra=2000 -> 0.197*2000^0.25 = 0.197*6.69=1.32 ; small mismatch at knee, acceptable schematic
\addplot[engred, very thick, domain=2000:1e6, samples=40] {0.197*(x^0.25)};
\draw[enggray, dotted] (axis cs:2000,0.8) -- (axis cs:2000,20);
\node[font=\scriptsize, enggray, anchor=south, rotate=90] at (axis cs:2000,1.3) {onset $\text{Ra}_c\approx 2000$};
\node[font=\scriptsize, engblue, anchor=south] at (axis cs:60,1.05) {still-air layer};
\node[font=\scriptsize, engred, anchor=south east] at (axis cs:6e5,10) {convection cells};
\end{axis}
\end{tikzpicture}
\caption[Nusselt number of an enclosed air gap]{An enclosed air gap does
not always convect. Below a critical gap Rayleigh number near $2000$ the
buoyant force cannot overcome viscous damping, no cells form, and the gap
transfers heat by pure conduction with a Nusselt number of one, the still
layer that gives a window air gap and a foam pore their insulating value.
Above the onset, convection rolls appear and the Nusselt number climbs as
roughly $\text{Ra}^{1/4}$, so a gap widened past the onset stops earning
its width because the new convection carries the extra heat the added
conduction resistance was meant to block. The design reading is that an
insulating gas gap is kept narrow enough to sit on the conduction floor,
which is why sealed-unit windows fix the gap near the width that keeps
$\text{Ra}$ below onset rather than making it as wide as possible.}
\label{fig:vol04ch05:enclosed-cavity}
\end{figure}
